\section{AHSME 1965}
        \begin{enumerate}
        	\item (\textit{AHSME 1965}) The number of real values of $x$ satisfying the equation $2^{2x^2 - 7x + 5} = 1$ is: 


 $\textbf{(A)}\ 0 \qquad  \textbf{(B) }\ 1 \qquad  \textbf{(C) }\ 2 \qquad  \textbf{(D) }\ 3 \qquad  \textbf{(E) }\ \text{more than 4}$


 

	\item (\textit{AHSME 1965}) A regular hexagon is inscribed in a circle. The ratio of the length of a side of the hexagon to the length of the 
shorter of the arcs intercepted by the side, is: 


 $\textbf{(A)}\ 1: 1 \qquad  \textbf{(B) }\ 1: 6 \qquad  \textbf{(C) }\ 1: \pi \qquad  \textbf{(D) }\ 3: \pi \qquad  \textbf{(E) }\ 6:\pi$


 

	\item (\textit{AHSME 1965}) The expression $(81)^{ - 2^{ - 2}}$ has the same value as: 


 $\textbf{(A)}\ \frac {1}{81} \qquad  \textbf{(B) }\ \frac {1}{3} \qquad  \textbf{(C) }\ 3 \qquad  \textbf{(D) }\ 81\qquad \textbf{(E) }\ 81^4$


 

	\item (\textit{AHSME 1965}) Line $\ell_2$ intersects line $\ell_1$ and line $\ell_3$ is parallel to $\ell_1$. The three lines are distinct and lie in a plane. 
The number of points equidistant from all three lines is: 


 $\textbf{(A)}\ 0 \qquad  \textbf{(B) }\ 1 \qquad  \textbf{(C) }\ 2 \qquad  \textbf{(D) }\ 4 \qquad  \textbf{(E) }\ 8$


 

	\item (\textit{AHSME 1965}) When the repeating decimal $0.363636\ldots$ is written in simplest fractional form, the sum of the numerator and denominator is: 


 $\textbf{(A)}\ 15 \qquad  \textbf{(B) }\ 45 \qquad  \textbf{(C) }\ 114 \qquad  \textbf{(D) }\ 135 \qquad  \textbf{(E) }\ 150$


 

	\item (\textit{AHSME 1965}) If $10^{\log_{10}9} = 8x + 5$ then $x$ equals: 


 $\textbf{(A)}\ 0 \qquad  \textbf{(B) }\ \frac {1}{2} \qquad  \textbf{(C) }\ \frac {5}{8} \qquad  \textbf{(D) }\ \frac{9}{8}\qquad \textbf{(E) }\ \frac{2\log_{10}3-5}{8}$


 

	\item (\textit{AHSME 1965}) The sum of the reciprocals of the roots of the equation $ax^2 + bx + c = 0$ is: 


 $\textbf{(A)}\ \frac {1}{a} + \frac {1}{b} \qquad  \textbf{(B) }\ - \frac {c}{b} \qquad  \textbf{(C) }\ \frac{b}{c}\qquad \textbf{(D) }\ -\frac{a}{b}\qquad \textbf{(E) }\ -\frac{b}{c}$


 

	\item (\textit{AHSME 1965}) One side of a given triangle is 18 inches. Inside the triangle a line segment is drawn parallel to this side forming a trapezoid 
whose area is one-third of that of the triangle. The length of this segment, in inches, is: 


 $\textbf{(A)}\ 6\sqrt {6} \qquad  \textbf{(B) }\ 9\sqrt {2} \qquad  \textbf{(C) }\ 12 \qquad  \textbf{(D) }\ 6\sqrt{3}\qquad \textbf{(E) }\ 9$


 

	\item (\textit{AHSME 1965}) The vertex of the parabola $y = x^2 - 8x + c$ will be a point on the $x$-axis if the value of $c$ is: 


 $\textbf{(A)}\ - 16 \qquad  \textbf{(B) }\ - 4 \qquad  \textbf{(C) }\ 4 \qquad  \textbf{(D) }\ 8 \qquad  \textbf{(E) }\ 16$


 

	\item (\textit{AHSME 1965}) The statement $x^2 - x - 6 < 0$ is equivalent to the statement: 


 $\textbf{(A)}\ - 2 < x < 3 \qquad  \textbf{(B) }\ x > - 2 \qquad  \textbf{(C) }\ x < 3 \\ \textbf{(D) }\ x > 3 \text{ and }x < - 2 \qquad  \textbf{(E) }\ x > 3 \text{ or }x < - 2$


 

	\item (\textit{AHSME 1965}) Consider the statements: $\newline$
$I: (\sqrt{-4})(\sqrt {-16}) = \sqrt{(-4)(-16)}$, $\newline$
$II: \sqrt{(-4)(-16)} = \sqrt{64}$, $\newline$
$III: \sqrt{64} = 8$. $\newline$
Of these the following are incorrect. 


 $\textbf{(A)}\ \text{none} \qquad  \textbf{(B) }\ \text{I only} \qquad  \textbf{(C) }\ \text{II only} \qquad  \textbf{(D) }\ \text{III only}\qquad \textbf{(E) }\ \text{I and III only}$


 

	\item (\textit{AHSME 1965}) A rhombus is inscribed in $\triangle ABC$ in such a way that one of its vertices is $A$ and two of its sides lie along $AB$ and $AC$. 
If $\overline{AC} = 6$ inches, $\overline{AB} = 12$ inches, and $\overline{BC} = 8$ inches, the side of the rhombus, in inches, is: 


 $\textbf{(A)}\ 2 \qquad  \textbf{(B) }\ 3 \qquad  \textbf{(C) }\ 3 \frac {1}{2} \qquad  \textbf{(D) }\ 4 \qquad  \textbf{(E) }\ 5$


 

	\item (\textit{AHSME 1965}) Let $n$ be the number of number-pairs $(x,y)$ which satisfy $5y - 3x = 15$ and $x^2 + y^2 \le 16$. Then $n$ is: 


 $\textbf{(A)}\ 0 \qquad  \textbf{(B) }\ 1 \qquad  \textbf{(C) }\ 2 \qquad  \textbf{(D) }\ \text{more than two, but finite}\qquad \textbf{(E) }\ \text{greater than any finite number}$


 

	\item (\textit{AHSME 1965}) The sum of the numerical coefficients in the complete expansion of $(x^2 - 2xy + y^2)^7$ is: 


 $\textbf{(A)}\ 0 \qquad  \textbf{(B) }\ 7 \qquad  \textbf{(C) }\ 14 \qquad  \textbf{(D) }\ 128 \qquad  \textbf{(E) }\ 128^2$


 

	\item (\textit{AHSME 1965}) The symbol $25_b$ represents a two-digit number in the base $b$. If the number $52_b$ is double the number $25_b$, then $b$ is: 


 $\textbf{(A)}\ 7 \qquad  \textbf{(B) }\ 8 \qquad  \textbf{(C) }\ 9 \qquad  \textbf{(D) }\ 11 \qquad  \textbf{(E) }\ 12$


 

	\item (\textit{AHSME 1965}) Let line $AC$ be perpendicular to line $CE$. Connect $A$ to $D$, the midpoint of $CE$, and connect $E$ to $B$, 
the midpoint of $AC$. If $AD$ and $EB$ intersect in point $F$, and $\overline{BC} = \overline{CD} = 15$ inches, 
then the area of triangle $DFE$, in square inches, is: 


 $\textbf{(A)}\ 50 \qquad  \textbf{(B) }\ 50\sqrt {2} \qquad  \textbf{(C) }\ 75 \qquad  \textbf{(D) }\ \frac{15}{2}\sqrt{105}\qquad \textbf{(E) }\ 100$


 

	\item (\textit{AHSME 1965}) Given the true statement: The picnic on Sunday will not be held only if the weather is not fair. We can then conclude that: 


 $\textbf{(A)}\ \text{If the picnic is held, Sunday's weather is undoubtedly fair.} \\ \textbf{(B) }\ \text{If the picnic is not held, Sunday's weather is possibly unfair.} \\ \textbf{(C) }\ \text{If it is not fair Sunday, the picnic will not be held.} \\ \textbf{(D) }\ \text{If it is fair Sunday, the picnic may be held.} \\ \textbf{(E) }\ \text{If it is fair Sunday, the picnic must be held.}$


 

	\item (\textit{AHSME 1965}) If $1 - y$ is used as an approximation to the value of $\frac {1}{1 + y}, |y| < 1$, the ratio of the error made to the correct value is: 


 $\textbf{(A)}\ y \qquad  \textbf{(B) }\ y^2 \qquad  \textbf{(C) }\ \frac {1}{1 + y} \qquad  \textbf{(D) }\ \frac{y}{1+y}\qquad \textbf{(E) }\ \frac{y^2}{1+y}\qquad$


 

	\item (\textit{AHSME 1965}) If $x^4 + 4x^3 + 6px^2 + 4qx + r$ is exactly divisible by $x^3 + 3x^2 + 9x + 3$, the value of $(p + q)r$ is: 


 $\textbf{(A)}\ - 18 \qquad  \textbf{(B) }\ 12 \qquad  \textbf{(C) }\ 15 \qquad  \textbf{(D) }\ 27 \qquad  \textbf{(E) }\ 45 \qquad$


 

	\item (\textit{AHSME 1965}) For every $n$ the sum of n terms of an arithmetic progression is $2n + 3n^2$. The $r$th term is: 


 $\textbf{(A)}\ 3r^2 \qquad  \textbf{(B) }\ 3r^2 + 2r \qquad  \textbf{(C) }\ 6r - 1 \qquad  \textbf{(D) }\ 5r + 5 \qquad  \textbf{(E) }\ 6r+2\qquad$


 

	\item (\textit{AHSME 1965}) It is possible to choose $x > \frac {2}{3}$ in such a way that the value of $\log_{10}(x^2 + 3) - 2 \log_{10}x$ is 


 $\textbf{(A)}\ \text{negative} \qquad  \textbf{(B) }\ \text{zero} \qquad  \textbf{(C) }\ \text{one} \\ \textbf{(D) }\ \text{smaller than any positive number that might be specified} \\ \textbf{(E) }\ \text{greater than any positive number that might be specified}$


 

	\item (\textit{AHSME 1965}) If $a_2 \neq 0$ and $r$ and $s$ are the roots of $a_0 + a_1x + a_2x^2 = 0$, then the equality 
$a_0 + a_1x + a_2x^2 = a_0\left (1 - \frac {x}{r} \right ) \left (1 - \frac {x}{s} \right )$ holds: 


 $\textbf{(A)}\ \text{for all values of }x, a_0\neq 0 \qquad \textbf{(B) }\ \text{for all values of }x \\ \textbf{(C) }\ \text{only when }x = 0 \qquad \textbf{(D) }\ \text{only when }x = r \text{ or }x = s \\ \textbf{(E) }\ \text{only when }x = r \text{ or }x = s, a_0 \neq 0$


 

	\item (\textit{AHSME 1965}) If we write $|x^2 - 4| < N$ for all $x$ such that $|x - 2| < 0.01$, the smallest value we can use for $N$ is: 


 $\textbf{(A)}\ .0301 \qquad  \textbf{(B) }\ .0349 \qquad  \textbf{(C) }\ .0399 \qquad  \textbf{(D) }\ .0401 \qquad  \textbf{(E) }\ .0499\qquad$


 

	\item (\textit{AHSME 1965}) Given the sequence $10^{\frac {1}{11}},10^{\frac {2}{11}},10^{\frac {3}{11}},\ldots,10^{\frac {n}{11}}$, 
the smallest value of n such that the product of the first $n$ members of this sequence exceeds $100000$ is: 


 $\textbf{(A)}\ 7 \qquad  \textbf{(B) }\ 8 \qquad  \textbf{(C) }\ 9 \qquad  \textbf{(D) }\ 10 \qquad  \textbf{(E) }\ 11$


 

	\item (\textit{AHSME 1965}) Let $ABCD$ be a quadrilateral with $AB$ extended to $E$ so that $\overline{AB} = \overline{BE}$. 
Lines $AC$ and $CE$ are drawn to form $\angle{ACE}$. For this angle to be a right angle it is necessary that quadrilateral $ABCD$ have: 


 $\textbf{(A)}\ \text{all angles equal} \qquad \textbf{(B) }\ \text{all sides equal} \\ \textbf{(C) }\ \text{two pairs of equal sides} \qquad \textbf{(D) }\ \text{one pair of equal sides} \\ \textbf{(E) }\ \text{one pair of equal angles}$


 

	\item (\textit{AHSME 1965}) For the numbers $a, b, c, d, e$ define $m$ to be the arithmetic mean of all five numbers; 
$k$ to be the arithmetic mean of $a$ and $b$; $l$ to be the arithmetic mean of $c, d$, and $e$; 
and $p$ to be the arithmetic mean of $k$ and $l$. Then, no matter how $a, b, c, d$, and $e$ are chosen, we shall always have: 


 $\textbf{(A)}\ m = p \qquad  \textbf{(B) }\ m \ge p \qquad  \textbf{(C) }\ m > p \qquad \\ \textbf{(D) }\ m < p\qquad \textbf{(E) }\ \text{none of these}$


 

	\item (\textit{AHSME 1965}) When $y^2 + my + 2$ is divided by $y - 1$ the quotient is $f(y)$ and the remainder is $R_1$. 
When $y^2 + my + 2$ is divided by $y + 1$ the quotient is $g(y)$ and the remainder is $R_2$. If $R_1 = R_2$ then $m$ is: 


 $\textbf{(A)}\ 0 \qquad  \textbf{(B) }\ 1 \qquad  \textbf{(C) }\ 2 \qquad  \textbf{(D) }\ - 1 \qquad  \textbf{(E) }\ \text{an undetermined constant}$


 

	\item (\textit{AHSME 1965}) An escalator (moving staircase) of $n$ uniform steps visible at all times descends at constant speed. 
Two boys, $A$ and $Z$, walk down the escalator steadily as it moves, A negotiating twice as many escalator 
steps per minute as $Z$. $A$ reaches the bottom after taking $27$ steps while $Z$ reaches the bottom after taking $18$ steps. Then $n$ is: 


 $\textbf{(A)}\ 63 \qquad  \textbf{(B) }\ 54 \qquad  \textbf{(C) }\ 45 \qquad  \textbf{(D) }\ 36 \qquad  \textbf{(E) }\ 30$


 

	\item (\textit{AHSME 1965}) Of $28$ students taking at least one subject the number taking Mathematics and English only equals the number 
taking Mathematics only. No student takes English only or History only, and six students take Mathematics and 
History, but not English. The number taking English and History only is five times the number taking all three subjects. 
If the number taking all three subjects is even and non-zero, the number taking English and Mathematics only is: 


 $\textbf{(A)}\ 5 \qquad  \textbf{(B) }\ 6 \qquad  \textbf{(C) }\ 7 \qquad  \textbf{(D) }\ 8 \qquad  \textbf{(E) }\ 9$


 

	\item (\textit{AHSME 1965}) Let $BC$ of right triangle $ABC$ be the diameter of a circle intersecting hypotenuse $AB$ in $D$. 
At $D$ a tangent is drawn cutting leg $CA$ in $F$. This information is not sufficient to prove that 


 $\textbf{(A)}\ DF \text{ bisects }CA \qquad  \textbf{(B) }\ DF \text{ bisects }\angle CDA \\ \textbf{(C) }\ DF = FA \qquad  \textbf{(D) }\ \angle A = \angle BCD \qquad  \textbf{(E) }\ \angle CFD = 2\angle A$


 

	\item (\textit{AHSME 1965}) The number of real values of $x$ satisfying the equality $(\log_ax)(\log_bx) = \log_ab$, where $a > 0, b > 0, a \neq 1, b \neq 1$, is: 


 $\textbf{(A)}\ 0 \qquad  \textbf{(B) }\ 1 \qquad  \textbf{(C) }\ 2 \qquad  \textbf{(D) }\ \text{a finite integer greater than 2}\qquad \textbf{(E) }\ \text{not finite}$


 

	\item (\textit{AHSME 1965}) An article costing $C$ dollars is sold for $100 at a loss of $x$ percent of the selling price. 
It is then resold at a profit of $x$ percent of the new selling price $S'$. 
If the difference between $S'$ and $C$ is $1\frac {1}{9}$ dollars, then $x$ is: 


 $\textbf{(A)}\ \text{undetermined} \qquad  \textbf{(B) }\ \frac {80}{9} \qquad  \textbf{(C) }\ 10 \qquad  \textbf{(D) }\ \frac{95}{9}\qquad \textbf{(E) }\ \frac{100}{9}$


 

	\item (\textit{AHSME 1965}) If the number $15!$, that is, $15 \cdot 14 \cdot 13 \dots 1$, ends with $k$ zeros when given to the base $12$ and ends with $h$ zeros 
when given to the base $10$, then $k + h$ equals: 


 $\textbf{(A)}\ 5 \qquad  \textbf{(B) }\ 6 \qquad  \textbf{(C) }\ 7 \qquad  \textbf{(D) }\ 8 \qquad  \textbf{(E) }\ 9$


 

	\item (\textit{AHSME 1965}) For $x \ge 0$ the smallest value of $\frac {4x^2 + 8x + 13}{6(1 + x)}$ is: 


 $\textbf{(A)}\ 1 \qquad  \textbf{(B) }\ 2 \qquad  \textbf{(C) }\ \frac {25}{12} \qquad  \textbf{(D) }\ \frac{13}{6}\qquad \textbf{(E) }\ \frac{34}{5}$


 

	\item (\textit{AHSME 1965}) The length of a rectangle is $5$ inches and its width is less than $4$ inches. The rectangle is folded so that two 
diagonally opposite vertices coincide. If the length of the crease is $\sqrt {6}$, then the width is: 


 $\textbf{(A)}\ \sqrt {2} \qquad  \textbf{(B) }\ \sqrt {3} \qquad  \textbf{(C) }\ 2 \qquad  \textbf{(D) }\ \sqrt{5}\qquad \textbf{(E) }\ \sqrt{\frac{11}{2}}$


 

	\item (\textit{AHSME 1965}) Given distinct straight lines $OA$ and $OB$. From a point in $OA$ a perpendicular is drawn to $OB$; 
from the foot of this perpendicular a line is drawn perpendicular to $OA$. 
From the foot of this second perpendicular a line is drawn perpendicular to $OB$; 
and so on indefinitely. The lengths of the first and second perpendiculars are $a$ and $b$, respectively. 
Then the sum of the lengths of the perpendiculars approaches a limit as the number of perpendiculars grows beyond all bounds. This limit is: 


 $\textbf{(A)}\ \frac {b}{a - b} \qquad  \textbf{(B) }\ \frac {a}{a - b} \qquad  \textbf{(C) }\ \frac {ab}{a - b} \qquad  \textbf{(D) }\ \frac{b^2}{a-b}\qquad \textbf{(E) }\ \frac{a^2}{a-b}$


 

	\item (\textit{AHSME 1965}) Point $E$ is selected on side $AB$ of $\triangle{ABC}$ in such a way that $AE: EB = 1: 3$ and point $D$ is selected on side $BC$ 
such that $CD: DB = 1: 2$. The point of intersection of $AD$ and $CE$ is $F$. Then $\frac {EF}{FC} + \frac {AF}{FD}$ is: 


 $\textbf{(A)}\ \frac {4}{5} \qquad  \textbf{(B) }\ \frac {5}{4} \qquad  \textbf{(C) }\ \frac {3}{2} \qquad  \textbf{(D) }\ 2\qquad \textbf{(E) }\ \frac{5}{2}$


 

	\item (\textit{AHSME 1965}) $A$ takes $m$ times as long to do a piece of work as $B$ and $C$ together; $B$ takes $n$ times as long as $C$ and $A$ together; 
and $C$ takes $x$ times as long as $A$ and $B$ together. Then $x$, in terms of $m$ and $n$, is: 


 $\textbf{(A)}\ \frac {2mn}{m + n} \qquad  \textbf{(B) }\ \frac {1}{2(m + n)} \qquad  \textbf{(C) }\ \frac{1}{m+n-mn}\qquad \textbf{(D) }\ \frac{1-mn}{m+n+2mn}\qquad \textbf{(E) }\ \frac{m+n+2}{mn-1}$


 

	\item (\textit{AHSME 1965}) A foreman noticed an inspector checking a $3$"-hole with a $2$"-plug and a $1$"-plug and suggested that two more gauges 
be inserted to be sure that the fit was snug. If the new gauges are alike, then the diameter, $d$, of each, to the nearest hundredth of an inch, is: 


 $\textbf{(A)}\ .87 \qquad  \textbf{(B) }\ .86 \qquad  \textbf{(C) }\ .83 \qquad  \textbf{(D) }\ .75 \qquad  \textbf{(E) }\ .71$


 

	\item (\textit{AHSME 1965}) Let $n$ be the number of integer values of $x$ such that $P = x^4 + 6x^3 + 11x^2 + 3x + 31$ is the square of an integer. Then $n$ is: 


 $\textbf{(A)}\ 4 \qquad  \textbf{(B) }\ 3 \qquad  \textbf{(C) }\ 2 \qquad  \textbf{(D) }\ 1 \qquad  \textbf{(E) }\ 0$


 

\end{enumerate}